% !PRINTED 103
\weberpage{103}

\noindent A function of the $n$th degree can be decomposed into $n$ linear factors.
\medskip

The magnitudes $\alpha,\beta,\ldots,\nu$ are then all roots of $f(x)$, and their number is therefore $n$. It is, however, not excluded that the same number occurs more than once among $\alpha,\beta,\ldots,\nu$. Then $f(x)$ would have fewer than $n$ roots, while the theorem would still remain that $f(x)$ is decomposable into $n$ linear factors. To establish agreement, it has been agreed that, if $x-\alpha$ goes into $f(x)$ several times, $\alpha$ is to be counted several times among the roots, and hence that one speaks of simple, double, triple, etc. roots.

By the concept of the derived functions [§ 13, (2)], if $\alpha$ is an arbitrary magnitude,
\begin{equation}
 f(x)=f(\alpha)+(x-\alpha)f'(\alpha)+\frac{(x-\alpha)^2}{1\cdot2}f''(\alpha)
 +\frac{(x-\alpha)^3}{1\cdot2\cdot3}f'''(\alpha)+\cdots .
\tag{4}
\end{equation}
If it is now again assumed that $f(\alpha)$ vanishes, then there results
\[
 f_1(x)=f'(\alpha)+\frac{x-\alpha}{1\cdot2}f''(\alpha)+\frac{(x-\alpha)^2}{1\cdot2\cdot3}f'''(\alpha)+\cdots,
\]
and hence
\[
 f_1(\alpha)=f'(\alpha).
\]
Thus $\alpha$ is a double root of $f(x)$ when $f'(\alpha)$ vanishes simultaneously with $f(\alpha)$. Formula (4) also shows that only under this assumption is $f(x)$ divisible by $(x-\alpha)^2$. This conclusion may be carried further and leads to the theorem:

\medskip
\noindent III. \gesperrt{The necessary and sufficient condition that $\alpha$ be an $m$-fold root of $f(x)$ is that $f(\alpha),f'(\alpha),f''(\alpha),\ldots,f^{(m-1)}(\alpha)$ vanish, while $f^{(m)}(\alpha)$ does not vanish.}

\begin{center}
{\large \S.\ 31.}

\vspace{0.5em}
{\large\gesperrt{Continuity of Integral Functions.}}
\end{center}

If, as before,
\[
 f(x)=a_0x^n+a_1x^{n-1}+\cdots+a_n
\]
is an integral function of $x$, we first prove the following theorem:

\medskip
